\chapter{研究计划}
\section{论文提纲}
计划中最终完成的学位论文将分为五章。全文以“注意力如何转化为可靠、可执行的定向交互”为主线，按照注意力由间接到直接、系统能力由目标选择到意图判别的顺序展开三项工作。各章内容安排如下：

第一章绪论：提出人机共生环境中的多对象定向交互问题，分析从目标发现、对象绑定、意图确认到会话维持与动作反馈的关键挑战；在此基础上，建立注意力外化、注意力对齐和意图判别逐级深入的研究框架，并概述本文的研究内容与主要贡献。

第二章基于手机指向的移动增强现实定向交互系统：介绍 RetroAR 的设计、实现与评估。该章研究如何将用户的手机指向转化为可计算的空间关注，利用可见光逆反射标签在同一光学链路中完成目标发现、身份读取与空间定位，并通过动态兴趣区、时空混合解码和六自由度位姿估计实现“所指即所控”。

第三章基于眼镜观察方向的原位双向交互系统：介绍 EchoSight 的设计、实现与评估。该章将定向交互入口由手持设备迁移至智能眼镜，研究双元件、双波段近红外光路和对齐检测机制，实现在无需预注册和额外无线配对条件下的原位双向交互，并分析由“所指即所控”向“所看即所控”推进所带来的自然性收益及其能力边界。

第四章眼、脑联合的注意力驱动意图交互系统：介绍 Omnisight 的设计、实现与评估。该章在 EchoSight 的基础上引入实测眼动，驱动可调光学发射单元主动对齐目标；进一步利用脑电辅助区分日常观看与主动操作，并研究定向光学选择通道与蓝牙持续会话通道之间的可验证绑定。该章将通过固定与可调光束对照、多模态消融和完整任务用户研究，检验“所思即所控”的可行性、实际收益与适用边界。

第五章总结与展望：总结手机指向、眼镜观察方向以及眼动与脑电三类注意力在对象选择、意图确认和持续交互中的递进关系，归纳注意力驱动定向交互的设计原则与适用条件，并讨论其在多对象人机共生环境中的后续研究方向。

\section{时间进度安排}
在前述三项工作中，RetroAR 与 EchoSight 的系统设计、原型实现和实验评估已经完成；Omnisight 的总体架构与实验问题已经明确，但眼镜原型、眼动与脑电同步、融合模型、混合链路协议及系统评估仍待完成。围绕 Omnisight 研究和学位论文撰写，计划设置以下四个时间节点：
\begin{enumerate}
  \item 在 2026 年 10 月 31 日以前，完成眼动追踪、可调光学发射和脑电采集模块的接入，建立统一时间戳与基础标定流程，完成眼动、头部姿态和光束坐标之间的转换，并记录系统同步误差与初步对齐误差。
  \item 在 2026 年 12 月 30 日以前，完成观看与主动操作任务范式、脑电预处理和离线分类基线，明确跨用户验证方法；同时完成光学目标标识与蓝牙设备身份的可验证绑定，以及目标切换、会话超时、重绑定和恢复机制，形成可执行的端到端 Omnisight 原型。
  \item 在 2027 年 2 月 31 日以前，完成固定与可调光束对照、眼动与脑电多模态消融及完整任务用户研究，系统评估目标错选率、误触发率、漏检率、决策与端到端时延、会话恢复时间、功耗和用户负担；在复核实验结果与系统边界的基础上，完成 Omnisight 研究论文的撰写与投稿。
  \item 在 2027 年 3 月 31 日以前，整合 RetroAR、EchoSight 与 Omnisight 三项工作，完成学位论文的撰写、修改和答辩准备。
\end{enumerate}
